\section{Formula derivations and implementation conventions}
\label{app:derivations}
Shared formulas are derived once, with explicit assumptions and published technical references.
Finite-array consequences are derived here; statistical methods and design
preferences are supplementary. Symbols are local: $N$ is a sampled record
length, $L$ the reference/reconstruction length, $K$ a transform length and
$L_q=2^B$ the number of quantiser levels. Powers specify their grid.
Receiver algorithms use no reference waveform or true channel errors.

\subsection{Finite-record Fourier, windows and power}
\label{app:dft}
\paragraph{DFT pair and one-sided scaling.}
Use the unscaled forward transform and $1/L$ inverse:
\begin{align}
 X[k]&=\sum_{n=0}^{L-1}x[n]e^{-j2\pi kn/L},\quad
 f_k=\frac{kF_0}{L},\quad\Delta f=\frac{F_0}{L}=\BinSpacing\,\mathrm{Hz},
 \label{eq:r11-dft}\\
 x[n]&=\frac1L\sum_{k=0}^{L-1}X[k]e^{j2\pi kn/L}.
 \label{eq:inverse}
\end{align}
Euler's cosine pair, $A\cos\theta=(A/2)(e^{j\theta}+e^{-j\theta})$,
gives, for an isolated coherent interior component,
\begin{equation}
 x[n]=A\cos(2\pi k_0n/L+\phi)
 \quad\Longrightarrow\quad X[k_0]=\tfrac{LA}{2}e^{j\phi}.
 \label{eq:cosine-amplitude}
\end{equation}
For $N$ actual samples and even $K\ge N$, the general windowed scales are

\begin{align}
 Y_{a,K}[k]&=\sum_{n=0}^{N-1}y[n]a[n]e^{-j2\pi kn/K},\quad
 f_k=\frac{kf_s}{K},\nonumber\\
 A_{a,K}[k]&=\frac{c_k|Y_{a,K}[k]|}{\sum_na[n]},\quad
 \widehat S_{a,K}[k]=\frac{c_k|Y_{a,K}[k]|^2}{f_s\sum_na[n]^2}.
 \label{eq:r22-spectrum-scales}
\end{align}
Here $c_k=1$ at DC/Nyquist and 2 at other positive bins. For $N=K=L$,
write $X_w=Y_{w,L}$ and $CG=L^{-1}\sum_nw[n]$; thus
$\widehat A_w=c_k|X_w|/(L\,CG)$. With $W=\operatorname{DFT}_L(w)$,
substitution of the cosine pair yields
\begin{equation}
 X_w[k_0]=\tfrac A2e^{j\phi}W[0]+\tfrac A2e^{-j\phi}W[2k_0],
 \qquad W[0]=L\,CG.
 \label{eq:window-gain-derivation}
\end{equation}
Normalising the first term does not remove mirror, off-bin or overlapping
components. Padding changes $\Delta f_K=f_s/K$, not the window-sum denominators.

\paragraph{Orthogonality, Parseval and PSD.}
For DFT indices $k,\ell\in\{0,\ldots,L-1\}$, the geometric ratio
$r=e^{j2\pi(k-\ell)/L}$ gives $(1-r^L)/(1-r)=0$ when $k\ne\ell$.
Insert the IDFT, cancel unequal-index terms, then pair conjugate bins:
\begin{align}
 \sum_{n=0}^{L-1}e^{j2\pi(k-\ell)n/L}&=L\delta_{k\ell},
 \label{eq:finite-orthogonality}\\
 \sum_n|x[n]|^2
 &=\frac1{L^2}\sum_{k,\ell}X[k]X^*[\ell]L\delta_{k\ell}
 =\frac1L\sum_k|X[k]|^2,\label{eq:parseval-steps}\\
 P_k&=c_k\frac{|X[k]|^2}{L^2},\qquad
 \widehat P=L^{-1}\sum_nx[n]^2=\sum_{k=0}^{L/2}P_k.
 \label{eq:parseval}
\end{align}
Applying the same identity to the padded windowed record gives
\begin{equation}
 \sum_{k=0}^{K/2}\widehat S_{a,K}[k]\frac{f_s}{K}
 =\frac{\sum_{k=0}^{K-1}|Y_{a,K}[k]|^2}{K\sum_na[n]^2}
 =\frac{\sum_ny[n]^2a[n]^2}{\sum_na[n]^2}.
 \label{eq:r22-psd-integral}
\end{equation}
PSD has squared NPY-amplitude/Hz units. Its weighted integral need not equal
the raw signal power used for noise generation.

\paragraph{Shared symmetric windows.}
For length $M$, $0\le n<M$, the tested windows are
\citep{numpy21hanning,numpy21hamming}
\begin{equation}
 \theta_n=\frac{2\pi n}{M-1},\quad w_{\rm rect}=1,\quad
 w_{\rm Hann}=\tfrac12(1-\cos\theta_n),\quad
 w_{\rm Hamming}=0.54-0.46\cos\theta_n.
 \label{eq:r22-windows}
\end{equation}
The spectral Hann $w_H$ uses $M=L$; kernel tapers use the FIR length.

\subsection{Sampling, interpolation and aliasing}
\label{app:sampling}
\paragraph{Folding and periodic recovery.}
For decimation $L=MN$, substitute $x_{\rm ref}[Mm]$ from the IDFT.
Orthogonality retains only $r=k+qN$, including phase:
\begin{equation}
 \frac{Y[k]}N=\frac1{LN}\sum_{r=0}^{L-1}X[r]
 \sum_{m=0}^{N-1}e^{j2\pi(r-k)m/N}
 =\sum_{q=0}^{M-1}\frac{X[k+qN]}L,\quad0\le k<N.
 \label{eq:decimation-folding}
\end{equation}
For real even-$N$ data, $Y=\operatorname{DFT}_N(x_M)$. Pairing the reference
IDFT and evaluating at $t/T=n/L$ gives the periodic interpolator:
\begin{equation}
 \widehat x[n]=\frac1N\left[Y[0]
 +2\operatorname{Re}\sum_{k=1}^{N/2-1}Y[k]e^{j2\pi kn/L}
 +Y[N/2]\cos\!\left(\frac{\pi Nn}{L}\right)\right].
 \label{eq:periodic-interpolation}
\end{equation}
For \texttt{irfft(Z,L)}, $L>N$, the old Nyquist coefficient splits into
an interior conjugate pair, requiring
\begin{equation}
 Z[k]=\frac LN Y[k]\ (k<N/2),\quad
 Z[N/2]=\frac L{2N}Y[N/2],\quad Z[k]=0\ (k>N/2).
 \label{eq:interpolation-scaling}
\end{equation}
For $L=N$ it is not split. DC/phased/Nyquist tests and independent SciPy
resampling check the convention \citep{numpy21irfft,scipy1153resample}.
The measured two-second periodicity supports this ideal-baseband recovery. Its inputs are only the samples and
target length; it does not recover arbitrary nonperiodic records.

\paragraph{Boundary phase and signed aliases.}
At $f_0=f_s/2$, the unknown sine component is lost:
\begin{equation}
 A\cos(\pi m+\phi)=A(-1)^m\cos\phi.
 \label{eq:critical-phase}
\end{equation}
For nonboundary $f_d=f_0-qf_s\in(-f_s/2,f_s/2)$, $f_d\ne0$, and $f_a=|f_d|$,
\begin{equation}
 A\cos(2\pi f_0m/f_s+\phi)=A\cos(2\pi f_am/f_s+\phi_a),\qquad
 \phi_a=\operatorname{sgn}(f_d)\phi.
 \label{eq:alias-phase}
\end{equation}
In the 8~kHz diagnostic, the retained low-frequency terms cancel, leaving
\begin{equation}
 e_8(t)\simeq-0.1\sin(2\pi\,2000t)-0.1\sin(2\pi\,6000t).
 \label{eq:alias-error-model}
\end{equation}
These distinct integer-bin tones are orthogonal on the common grid, so
\begin{equation}
 \mathrm{MSE}_{8,\rm pred}=0.1^2/2+0.1^2/2=0.01.
 \label{eq:alias-error-power}
\end{equation}
This explains measured phases after the experiment; no source tones enter
the recovery algorithm.

\subsection{Linear filtering, finite boundaries and reconstruction energy}
\label{app:filter}
\paragraph{Convolution and padding.}
The reference impulse expansion $y[n]=\sum_my[m]\delta[n-m]$, linearity
and time invariance give 
\begin{equation}
 z[n]=\mathcal T\{y\}[n]=\sum_my[m]\mathcal T\{\delta[\cdot-m]\}[n]
 =\sum_my[m]h[n-m].
 \label{eq:r21-lti-convolution}
\end{equation}
For circular convolution, $r=(n-m)\bmod K$ separates the DFT sum:
\begin{equation}
 Z_c[k]=\sum_mh[m]e^{-j2\pi km/K}\sum_ry[r]e^{-j2\pi kr/K}
 =H_K[k]Y_K[k].
 \label{eq:r21-convolution-theorem}
\end{equation}
Zero extension gives $N+M-1$ full outputs. Linear convolution is therefore
\begin{equation}
 z_K[n]=\frac1K\sum_{k=0}^{K-1}H_K[k]Y_K[k]e^{j2\pi kn/K},
 \qquad K\ge N+M-1.
 \label{eq:r21-padded-convolution}
\end{equation}
retaining the first $N+M-1$ values, with zero padding.
The 513-tap prototype has 32512 full outputs; $K=32768$ is convenient,
not a NumPy requirement. An insufficient-padding counterexample is

\begin{equation}
 [1,2,1]*[2,0,2]=[2,4,4,4,2],
 \qquad
 \operatorname{IDFT}_4(Y_4H_4)=[4,4,4,4].
 \label{eq:r21-wrap-example}
\end{equation}
A centred three-point mean has aligned response
\begin{equation}
 h=[1,1,1]/3,\quad d=1,\qquad
 G_3(\omega)=e^{j\omega d}\sum_{m=0}^2h[m]e^{-j\omega m}
 =\frac{1+2\cos\omega}{3}.
 \label{eq:r21-ma-response}
\end{equation}

\paragraph{One finite-IDFT filter family.}
For signed frequencies $f_k$, measured corrupted-spectrum centres $f_i$,
half-width $b$, odd $M$ and $d=(M-1)/2$,

\begin{align}
 B_b[k]&=\mathbf1\!\left\{\min_i\bigl||f_k|-f_i\bigr|\le b\right\},\\
 h_{M,b,w}[m]&=w[m]\frac1N\sum_{k=0}^{N-1}
 B_b[k]e^{j2\pi k(m-d)/N},\qquad d=(M-1)/2.
 \label{eq:r22-family}
\end{align}
Use the symmetric windows in Eq.~\eqref{eq:r22-windows} on the kernel,
$0\le m<M$. The Hann prototypes use 513 taps and 80~Hz half-width in R2.1,
then 385 taps and 70~Hz in R2.2. The low-pass comparator uses
$B_b=\mathbf1_{\{|f_k|\le6600\,\mathrm{Hz}\}}$.
This project construction uses the declared R1 priors, not true noise.
The realised response is $\operatorname{DFT}_K(h)$, not the binary template.

With zero extension, the reference-independent crop is
\begin{equation}
 \widehat x[n]=z_{\rm full}[n+d],\quad 0\le n<N.
 \label{eq:r21-centred-crop}
\end{equation}
The offline crop uses future samples; causal FIR delay is $d/f_s$
(16~ms for the prototype), distinct from whole-record FFT buffering.

\paragraph{Finite-boundary noise gain, distortion and crossover.}
For a fixed filter, write $d_x=F(x)-x$, $v=F(e)$ and $y=x+e$:
\begin{equation}
 F(y)-x=d_x+v,\qquad
 \mathrm{MSE}=(\|d_x\|^2+\|v\|^2+2d_x^{\mathsf T}v)/N.
 \label{eq:r22-distortion-decomposition}
\end{equation}
For new independent zero-mean noise, $g[\ell]=h[\ell+d]$, $d<N$,
each lag overlaps $N-|\ell|$ positions. Thus
\begin{align}
 G_{\rm finite}&=\sum_{\ell=-d}^dg[\ell]^2(1-|\ell|/N),\nonumber\\
 \mathbb E[\mathrm{MSE}]&=D+\sigma^2G_{\rm finite},\qquad D=\|d_x\|^2/N.
 \label{eq:r22-finite-noise-gain}
\end{align}
The cross term vanishes in expectation, not in each trial. For $D>0,G<1$,
$D+\sigma^2G=\sigma^2$ gives $\sigma^2=D/(1-G)$ and
\begin{equation}
 S_* =10\log_{10}\frac{P_x(1-G_{\rm finite})}{D}
       \approx\RttCrossover\ {\rm dB}.
 \label{eq:r22-crossover}
\end{equation}

\paragraph{Reconstruction-grid energy and expected noise.}
For the frozen filter, $D=\RttCleanMse$ and $G=\RttNoiseGain$ in
Eq.~\eqref{eq:r22-finite-noise-gain}. Let $V=\operatorname{DFT}_N(v)$ for
a real even-$N$ error vector and $A_\pi=V[N/2]/N=N^{-1}\sum_n(-1)^nv[n]$.
For interpolation $R$ to $L=3N$, ordinary pairs retain their mean-square
contribution; splitting the old Nyquist coefficient halves its power:
\begin{align}
 \frac1N\|v\|^2
 &=\frac{|V[0]|^2+
 2\sum_{k=1}^{N/2-1}|V[k]|^2+|V[N/2]|^2}{N^2},\nonumber\\
 \frac1L\|R(v)\|^2
 &=\frac{|V[0]|^2+
 2\sum_{k=1}^{N/2-1}|V[k]|^2+\tfrac12|V[N/2]|^2}{N^2}\nonumber\\
 &=\frac1N\|v\|^2-\frac12 A_\pi^2.
 \label{eq:r23-reconstruction-energy}
\end{align}
With $a_n=(-1)^n/N$, independence gives
$\mathbb E[(a^{\mathsf T}Fe)^2]=\sigma^2\|F^{\mathsf T}a\|^2$, so
\begin{equation}
 G_{48}=G-\frac12\|F^\mathsf{T}a\|^2
       =\RtvNoiseGainHigh,\qquad
 G_{0,48}=1-\frac1{2N}=0.999984375.
 \label{eq:r23-high-gain}
\end{equation}
The 48~kHz predictions are $D_{48}+\sigma^2G_{48}$ and
$\sigma^2G_{0,48}$, with measured clean distortion $D_{48}=\RtvCleanHigh$;
negligible clean interpolation residuals are omitted. Actual errors use
the original reference, rather than assuming the two grids are equivalent.

\paragraph{Independent finite cosine-series kernel.}
For disjoint interior bands centred on DFT bins, integer $J_b=bN/f_s$
and $C_b=2J_b+1$, the finite geometric series gives an independent check:
\begin{align}
 S_{C_b}(\ell)
 &=\sum_{q=-J_b}^{J_b}e^{j2\pi q\ell/N}
 =\begin{cases}
 C_b,&\ell=0,\\
 \dfrac{\sin(C_b\pi\ell/N)}{\sin(\pi\ell/N)},&\ell\ne0,
 \end{cases}\nonumber\\
 h[m]&=\frac{2w[m]}{N}S_{C_b}(m-d)
       \sum_i\cos\!\left(\frac{2\pi f_i(m-d)}{f_s}\right).
 \label{eq:r23-audit-kernel}
\end{align}

\subsection{Quantisation, binary labels and amplitude errors}
\label{app:quant}
\paragraph{Endpoint codebook and shifted rounding.}
The endpoint mapping, Eq.~\eqref{eq:r31mapping}, has spacing
$\Delta=2A/(L_q-1)$ versus the ideal-model $2A/L_q$. The normalised epsilon
shifts rounding by $A\epsilon$, giving $|z_q-z_c|\le\Delta/2+A\epsilon$
apart from floating-point effects. The even endpoint codebook omits zero;
the retained shift maps zero to $-\Delta/2$. With clipping,
$e_q=(z_c-z)+(z_q-z_c)$, whose power includes their cross term.
Actual errors need not be white, independent or zero-mean
\citep{widrow1996quantization}.

\paragraph{Common headroom and an overload bound.}
Both branches share a range fixed before noise generation for the
$-10$ to 40~dB stress envelope. Let $M_b=\max_n|F_b(x)[n]|$,
$g_0=1$, $g_F=\|h\|_2$ and
$\sigma_{\max}=\sqrt{P_x10^{-S_{\min}/10}}=2.701851$ at $S_{\min}=-10$~dB.
Independence gives $\operatorname{Var}[(F\eta)[n]]
\le\sigma^2\|h\|_2^2$, including boundaries. Six-sigma headroom yields
\begin{equation}
 A=\left\lceil\max_b(M_b+6\sigma_{\max}g_b)\right\rceil
 =\left\lceil\max(\QRawHeadroom,\QFilterHeadroom)\right\rceil=\QHalfRange.
 \label{eq:r31range}
\end{equation}
For unbounded Gaussian noise this is not a zero-overload guarantee:
$\Pr(\text{any overload})\le2N\operatorname{erfc}(6/\sqrt2)=\QUnionBound$
for the two $N=32000$ branches by the union bound, without branch
independence. Overloads are counted before clipping; the stress envelope
does not expand R2's design claims.

\paragraph{Uniform-error integral and SQNR conventions.}
The reference's unoverloaded model has $p_e(e)=1/\Delta$ on
$[-\Delta/2,\Delta/2)$ and $\mathbb E[e]=0$ by symmetry. Its model variance is
\citep{bennett1948spectra}
\begin{equation}
 \mathbb E[e^2]=\frac1\Delta\int_{-\Delta/2}^{\Delta/2}e^2\,de
 =\frac1\Delta\left[\frac{e^3}3\right]_{-\Delta/2}^{\Delta/2}
 =\frac{\Delta^2}{12}.
 \label{eq:r31variance}
\end{equation}
For endpoint spacing and $R_a=2A$, this gives the model
\begin{equation}
 {\rm SQNR}_{\rm model}
 =10\log_{10}\!\left(\frac{12P_z}{\Delta^2}\right)
 =10\log_{10}\!\left(\frac{12P_z(2^B-1)^2}{R_a^2}\right).
 \label{eq:r31model}
\end{equation}
A full-range sine has $P_z=R_a^2/8$, so the ideal uniform spacing gives
$10\log_{10}(1.5\,4^B)=10\log_{10}1.5+20B\log_{10}2
\simeq1.76+6.02B$~dB. Endpoint spacing adds $20\log_{10}(1-2^{-B})$.
Neither sine formula is the actual multitone/noisy-input SQNR.

\paragraph{Random-label power and the p=0.5 full-grid check.}
At $p=0.5$, received bits are independent fair bits, independent of the
payload. Hamming(7,4)'s 16 equal eight-word correction spheres give uniform
decoded messages, hence independent fair information bits and uniform
disjoint 14-bit labels. For uniform $j\in\{0,\ldots,L_q-1\}$,
$\mathbb E[j]=(L_q-1)/2$ and $\operatorname{Var}(j)=(L_q^2-1)/12$, giving
\begin{equation}
 V_q=\left(\frac{2A}{L_q-1}\right)^2\operatorname{Var}(j)
 =\frac{A^2(L_q+1)}{3(L_q-1)}.
 \label{eq:r43-uniform-label}
\end{equation}
Here $A=18,L_q=2^{14}$. This is decoded-level power, not $\Delta^2/12$.
The reference cross term has zero expectation; Nyquist correction gives
\begin{equation}
 \mathbb E[d\mid p=0.5]
 =1+\frac{V_q}{P_{\rm ref}}\left(1-\frac1{2N}\right)
 =\WHalfTheory,\qquad N=32000.
 \label{eq:r43-half-theory}
\end{equation}

\paragraph{Bit depth, label significance and code-block coupling.}
The endpoint model predicts only the quantisation-component ratio
\begin{equation}
 \frac{P_{q,16}}{P_{q,14}}
 =\left(\frac{2^{14}-1}{2^{16}-1}\right)^2\simeq\frac1{16},
 \label{eq:r51-quant-ratio}
\end{equation}
For true MSB-first bits $b_{n,u}$ and error indicators $w_{n,u}$, XOR is
$\widehat b_{n,u}=b_{n,u}+(1-2b_{n,u})w_{n,u}$. Therefore
\begin{equation}
 \begin{aligned}
 \ell[n]&=\sum_{u=0}^{B-1}2^{B-1-u}b_{n,u},\quad z_q[n]=-A+\Delta_B\ell[n],\\
 \delta z[n]&=\Delta_B(\widehat\ell[n]-\ell[n])
 =\Delta_B\sum_{u=0}^{B-1}(1-2b_{n,u})2^{B-1-u}w_{n,u}.
 \end{aligned}
 \label{eq:r51-bit-coupling}
\end{equation}
At $B=14$, one MSB error has $2^{13}$ times an LSB error's amplitude
magnitude and $2^{26}$ times its squared magnitude. Squaring retains joint
$w_{n,u}w_{n,v}$ terms. Label starts $(Bn)\bmod4$ alternate 0/2 for 14 bits
but stay 0 for 16; dependent within-codeword errors therefore interact
differently with significant bits. Component/cross-term measurements capture
the combined change, not an isolated alignment effect or resolution-only gain.

\paragraph{Fixed-resolution packing diagnostic.}
With fixed 14-bit labels, $\alpha_{n,u}=\Delta 2^{13-u}(1-2b_{n,u})$.
Let $A_b$ place codeword $b$'s four slopes in their sample positions:
$\delta z=\sum_bA_bw_b$. The 128 BSC patterns give residual-bit mean $q$
and covariance $C$. Independent blocks give
$m=\mathbb E[\delta z]=\sum_bA_bq$ and
$\Sigma=\operatorname{Cov}(\delta z)=\sum_bA_bCA_b^T$.
Apply Eq.~\eqref{eq:r23-reconstruction-energy} to the zero-mean fluctuation:
\begin{equation}
 \mathbb E[d]=\frac{1}{P_{\rm ref}}\left\{
 \frac{\|\mathcal R(z_q+m)-x_{\rm ref}\|^2}{L}
 +\frac{\operatorname{tr}\Sigma}{N}
 -\frac{a^T\Sigma a}{2N^2}\right\},\quad a_n=(-1)^n.
 \label{eq:packing-expectation}
\end{equation}
The final term corrects split Nyquist energy. Packing changes $A_b$,
with $q,C,z_q$ fixed; Gaussian is absent. The code accumulates block
moments without storing $\Sigma$.

\subsection{Systematic ECC, channel error laws and payload counting}
\label{app:ecc}
Counting fixed-width samples gives
\begin{equation}
 R_{\mathrm{raw}}=Bf_s,\qquad N_{\mathrm{bits}}=BN,\qquad
 \frac{R_{\mathrm{raw},i}}{R_{\mathrm{raw},16k}}=\frac{f_{s,i}}{16000}.
 \label{eq:sampling-payload}
\end{equation}

\paragraph{Fixed matrices, syndrome, distance and miscorrection.}
The reference distance and correction radius are

\begin{align}
 d_H(\mathbf u,\mathbf v)&=\sum_{i=0}^{n-1}
       \mathbb{1}\{u_i\ne v_i\},&
 d_{\min}&=\min_{\mathbf u\ne\mathbf v\in C}d_H(\mathbf u,\mathbf v),
 \label{eq:r32-distance}\\
 R&=\frac{k}{n},&
 t&=\left\lfloor\frac{d_{\min}-1}{2}\right\rfloor.
 \label{eq:r32-rate}
\end{align}
A validity check detects up to $d_{\min}-1$ errors; correction guarantees
only $t$. All arithmetic is in $\mathrm{GF}(2)$, using integer NumPy
matrix multiplication modulo two over GF(2). The actual row-message matrices are

\begin{equation}
 G_7=\left[\begin{array}{rrrr|rrr}
 1&0&0&0&1&1&0\\0&1&0&0&1&0&1\\
 0&0&1&0&0&1&1\\0&0&0&1&1&1&1
 \end{array}\right],\qquad
 H_7=\left[\begin{array}{rrrr|rrr}
 1&1&0&1&1&0&0\\1&0&1&1&0&1&0\\0&1&1&1&0&0&1
 \end{array}\right].
 \label{eq:r32-matrices}
\end{equation}
For $G=[I_k\mid P]$ and $H=[P^T\mid I_{n-k}]$,
\begin{equation}
 GH^T=P+P=0,\quad \mathbf c=\mathbf mG,\quad
 \mathbf r=\mathbf c+\mathbf e,\quad
 \mathbf s=\mathbf rH^T=\mathbf eH^T.
 \label{eq:r32-syndrome}
\end{equation}
A single-error syndrome equals the corresponding column of $H$.
Its columns, in order, are 110, 101, 011, 111, 100, 010, 001.
Use column matching rather than binary-syndrome indexing;
zero syndrome leaves the word unchanged. Extract the first $k$ bits
after correction.
The reference example $1100\mapsto1100011$ becomes $1110011$ after flipping
zero-based position 2; syndrome $011$ restores it. Enumeration of all 16
messages gives minimum nonzero weight 3. Linearity,
$d_H(c_1,c_2)=w_H(c_1+c_2)$, proves $d_{\min}=3$.
The disjoint radius-one spheres contain $2^4(1+7)=2^7$ words, a perfect
partition; sphere size alone is not proof of correction.

Two flips $0000000\to1100000$ also give syndrome $011$: correction produces
the wrong valid word $1110000$, message $1110$, with final syndrome zero.
This decoder is not SECDED; validity does not prove message recovery.

\paragraph{Alternative codes and majority decoding.}
For the $(15,11)$ reference code, define $G_{15}=[I_{11}\mid P_{15}]$ and
$H_{15}=[P_{15}^T\mid I_4]$ by these ordered four-bit parity rows:
\begin{equation}
 \begin{aligned}
 \operatorname{rows}(P_{15})={}&(1100,1010,1001,0110,0101,0011,\\
 &1110,1101,1011,0111,1111).
 \end{aligned}
 \label{eq:r32-p15}
\end{equation}
Each string is one binary row. Distinct nonzero columns of $H_{15}$ exclude
weight-one/two codewords; the first generator row has weight three.
Thus $d_{\min}=3,t=1,R=11/15$. The five-repetition reference alternative has
$G_5=[1\ 1\ 1\ 1\ 1]$, $H_5=[\mathbf1_4\mid I_4]$,
$d_{\min}=5,t=2,R=1/5$. Majority fails on three or more flips:
\begin{equation}
 P_{b,5}=10p^3(1-p)^2+5p^4(1-p)+p^5.
 \label{eq:r32-repeat}
\end{equation}
The ideal-model three-repetition example instead gives $3p^2-2p^3$. Uncoded $(1,1)$ has
$d_{\min}=1,t=0,R=1,P_b=p$.

\paragraph{Exact framing, padding and effective rate.}
For $N$ labels of $B$ bits, duration $T$ and an $(n,k)$ code,
\begin{align}
 M_b&=NB,\qquad a=(-M_b)\bmod k,\qquad
 N_c=\left\lceil M_b/k\right\rceil,\nonumber\\
 M_t&=nN_c,\qquad R_{\mathrm{eff}}=M_b/M_t,\qquad r_t=M_t/T.
 \label{eq:r32-framing}
\end{align}
Here $N=32000,B=14,T=2$~s, so $M_b=448000$. R3.2 also uses $M=M_b,v_t=r_t$.
The last word has $b=M_b-k(N_c-1)$ real bits. Padding is encoded and counted
in the logical sending rate but excluded from real-payload BER; headers
and physical energy lie outside this rate definition.

Independent BSC flips imply
\begin{equation}
 E_i\overset{\mathrm{iid}}{\sim}\operatorname{Bernoulli}(p),\quad
 J_b=\sum_{i=1}^n E_i\sim\operatorname{Binomial}(n,p),\quad
 \Pr(J_b=j)=\binom njp^j(1-p)^{n-j}.
 \label{eq:r32-binomial}
\end{equation}

\paragraph{Exact complete-block failure and decoded BER.}
For the verified perfect Hamming and odd-repetition decoders, complete
padded-message failure occurs exactly when error weight exceeds $t$:
\begin{equation}
 q_{n,t}(p)=1-\sum_{j=0}^t\binom njp^j(1-p)^{n-j}
 =\sum_{j=t+1}^n\binom njp^j(1-p)^{n-j}.
 \label{eq:r32-bler}
\end{equation}
The positive tail is numerically stable. This law is not assumed for
arbitrary ECC; real-payload BER excludes Hamming(15,11)'s eight padding bits.
Translation invariance, Eq.~\eqref{eq:r42-pair-invariant}, permits all-zero
message enumeration. Grouping patterns by weight gives
\begin{equation}
 \begin{aligned}
 A_{c,j}&=\sum_{\mathbf e:\,w_H(\mathbf e)=j}w_H(\mathcal D_c(\mathbf e)),\\
 P_{b,c}(p)&=\frac1k\sum_{j=0}^n A_{c,j}p^j(1-p)^{n-j}.
 \end{aligned}
 \label{eq:r32-exactber}
\end{equation}
Multiplicity is included in $A_{c,j}$. The (7,4) coefficients
$(0,0,36,76,64,48,28,4)$ give $P_b=9p^2+O(p^3)$; (15,11) gives
$21p^2+O(p^3)$. Also $q_{7,1}=1-(1-p)^7-7p(1-p)^6=21p^2+O(p^3)$.

\paragraph{Capacity and rate are distinct from finite-code error.}
The memoryless BSC capacity is
\begin{equation}
 H_2(p)=-p\log_2p-(1-p)\log_2(1-p),\qquad C(p)=1-H_2(p),
 \label{eq:r33-capacity}
\end{equation}
Use $0\log_20=0$; $C(p)=4/7$ gives $p=\PCapacityRoot$.
Capacity does not predict this finite decoder's BER.

\paragraph{Exact finite-code BER crossover.}
Expanding the (7,4) enumerator and factoring gives the nontrivial
crossover in $(0,1/2)$:
\begin{align}
 P_{b,7}(p)&=9p^2-26p^3+30p^4-12p^5,\nonumber\\
 P_{b,7}(p)-p&=-p(2p-1)(p-1)(6p^2-6p+1),\label{eq:r33-bercross}\\
 p_{\rm BER}&=(3-\sqrt3)/6=\PBerRoot.\label{eq:r33-berroot}
\end{align}

\paragraph{Decoder invariance under a shared mask.}
Adding a valid codeword leaves syndrome and correction vector
$\kappa(\mathbf s)$ unchanged. For information projection $\pi_4$,
\begin{equation}
 \mathbf s=(\mathbf c\oplus\mathbf e)H^T=\mathbf eH^T,\qquad
 D_1(E_1(\mathbf b)\oplus\mathbf e)
 =\mathbf b\oplus\pi_4[\mathbf e\oplus\kappa(\mathbf eH^T)].
 \label{eq:r42-pair-invariant}
\end{equation}
Hence $\mathcal D(c\oplus e)=m\oplus\mathcal D(e)$; odd-repetition majority
has the same translation property. The same mask gives identical decoded
error patterns across messages.

\paragraph{Exact real-payload probability with a partial final word.}
Let $b=M_b-(N_c-1)k$ and let $a_{b,w}$ count weight-$w$ patterns whose
decoded first $b$ bits remain correct. Translation invariance and independent
disjoint BSC words give
\begin{align}
 s_b(p)&=\sum_{w=0}^n a_{b,w}p^w(1-p)^{n-w},\label{eq:r52-prefix}\\
 P_{\rm exact\ payload}
 &=(1-q_{n,t}(p))^{N_c-1}s_b(p).\label{eq:r52-payload}
\end{align}
For (15,11), $b=3$ and eight bits are padding, so
$(1-q_{15,1})^{40728}$ is not exact real-payload recovery.
For unpadded (7,4), $b=4$, $s_b=1-q_{7,1}$, yielding
$(1-q_{7,1})^{112000}$; uncoded transmission gives $(1-p)^{448000}$.

\subsection{Reconstruction metrics and error accounting}
\label{app:metrics}
\paragraph{Common finite-record metrics and units.}
For reference $u$, recovery $\widehat u$, $Q$ grid positions and $P_u>0$,
\begin{equation}
 \begin{aligned}
 a[j]&=\widehat u[j]-u[j],\quad P_u=Q^{-1}\sum_ju[j]^2,\quad
 E=\operatorname{MSE}=Q^{-1}\sum_ja[j]^2,\\
 \operatorname{RMSE}&=\sqrt E,\qquad \operatorname{NMSE}=E/P_u,\qquad
 S_{\rm error}=10\log_{10}(P_u/E).
 \end{aligned}
 \label{eq:r23-metrics}
\end{equation}

Cauchy--Schwarz gives
\begin{equation}
 \operatorname{MAE}=Q^{-1}\sum_j|a[j]|\leq
 \sqrt{Q^{-1}\sum_ja[j]^2}=\operatorname{RMSE}.
 \label{eq:r41-mae}
\end{equation}
MAE is not a monotone transform of MSE.

For paired repeat $r$,
\begin{equation}
 \delta_r=E_{0,r}-E_{F,r},\qquad
 g_r=S_{{\rm error},F,r}-S_{{\rm error},0,r}
     =10\log_{10}\frac{E_{0,r}}{E_{F,r}}.
 \label{eq:r23-paired}
\end{equation}

\paragraph{Boundary diagnostics and the aggregate denominator.}
For edge set $B$ of size $2q$ and complement $I$,
\begin{align}
 E_B&=\frac1{2q}\sum_{j\in B}a[j]^2,&
 E_I&=\frac1{Q-2q}\sum_{j\in I}a[j]^2,\nonumber\\
 E&=\frac{2q}{Q}E_B+\frac{Q-2q}{Q}E_I,&
 \eta_B&=\frac{(2q/Q)\overline E_B}{\overline E}.
 \label{eq:r23-boundary-metrics}
\end{align}
The aggregate ratio is not the mean of individual edge/total ratios.

\paragraph{Component budget and total error.}
The squared-sum identity in Eq.~\eqref{eq:r42-component-power} applies
also to upstream error $a$ and extra quantisation $q$. Define
$\rho_q=\operatorname{mean}(q^2)/\operatorname{mean}(a^2)$ for nonzero
upstream power. Cauchy--Schwarz bounds their cross term, giving
\begin{equation}
 \left|\frac{\operatorname{mean}[(a+q)^2]}
                  {\operatorname{mean}(a^2)}-1\right|
 \le \rho_q+2\sqrt{\rho_q}.
 \label{eq:r31-total-bound}
\end{equation}
A 1\% component budget permits a conservative 21\% total-MSE change;
the observed largest increase is \QMaxTotalPercent\%.

Expanding centred squares with $\bar e=L^{-1}\sum_\ell e[\ell]$ gives
\begin{equation}
 V_e=L^{-1}\sum_\ell(e[\ell]-\bar e)^2=D-\bar e^2,
 \qquad D=V_e+\bar e^2.
 \label{eq:r41-bias}
\end{equation}
Variance alone omits bias. Waveform MSE uses $L$, not the $L-1$ sample-variance divisor.

For centred reference/recovery arrays $x_c,y_c$,
\begin{equation}
 \rho=\frac{\sum_\ell x_c[\ell]y_c[\ell]}
 {\sqrt{\sum_\ell x_c[\ell]^2\sum_\ell y_c[\ell]^2}}.
 \label{eq:r41-correlation}
\end{equation}

\paragraph{Five components and ten signed cross terms.}
The five full-grid errors are
\begin{align}
 e_s&=\mathcal R x-x_{\rm ref},&
 e_f&=\mathcal R F_fx-\mathcal R x,&
 e_g&=\mathcal R F_f\eta,\nonumber\\
 e_q&=\mathcal R z_q-\mathcal R z_f,&
 e_b&=\widehat x-\mathcal R z_q,
 \label{eq:r42-components}
\end{align}
Expanding their squared sum without independence assumptions gives
\begin{equation}
 \widehat x-x_{\rm ref}=\sum_i e_i,\qquad
 D=\sum_i\frac{\|e_i\|^2}{L}
   +2\sum_{i<j}\frac{\langle e_i,e_j\rangle}{L}.
 \label{eq:r42-component-power}
\end{equation}
All five powers and ten signed cross terms are saved; they are not
normalised into a 100\% responsibility chart.

For repeated complete-record NMSE values,
\begin{equation}
 \bar d=R^{-1}\sum_rd_r,\qquad
 S_{\rm err,from\ mean}=-10\log_{10}\bar d.
 \label{eq:r41-average}
\end{equation}
Generally this differs from averaging per-record dB values. Exact zero
error gives $+\infty$; no epsilon floor is inserted.

\subsection{Uncertainty, planning and decision rules}
\label{app:statistics}
\paragraph{One mean/variance/interval convention.}
For $R$ independent complete-record/bundle values $v_r$,
\begin{align}
 \bar v&=R^{-1}\sum_{r=0}^{R-1}v_r,&
 s_v^2&=(R-1)^{-1}\sum_r(v_r-\bar v)^2,\\
 [L_v,U_v]&=\bar v\pm t_{1-\alpha/(2m),R-1}\frac{s_v}{\sqrt R}.
 \label{eq:app-mean-interval}
\end{align}
Here $\alpha=0.05$. Independence gives
$\operatorname{Var}(\bar v)=\operatorname{Var}(v_r)/R$.
These are approximate t/CLT mean intervals for potentially nonnormal
records. Bonferroni's union bound
$\Pr(\bigcap_{j=1}^mA_j)\ge1-\sum_j\Pr(A_j^c)$ needs no independence
between statistics, but does not make marginal coverage exact
\citep{nistmeaninterval,nistbonferroni,dunn1961comparisons,scipy1153t,scipy1153sem}.
Applications use $(R,m)=(50,1)$ in R2.3, $(5,1)$ in R3.2,
$(300,30)$ in R3.4, $(1400,92)$ in R4.3 and $(350,19)$ in R5.
Pilot mean/SD estimates use 20 bundles and never enter confirmation means.
Paired differences and objectives are formed within bundles first;
correlated time positions and bits do not increase $R$.
These statistical methods supplement the DSP reference.

For $M_c$ independent disjoint codewords with failure probability $q(p)$,
\begin{equation}
 F\sim\operatorname{Binomial}(M_c,q(p)),\qquad
 \mathbb E[F]=M_cq(p),\qquad
 \Pr(F=0)=(1-q(p))^{M_c}.
 \label{eq:r33-exposure}
\end{equation}

\paragraph{Exact proportion intervals and the zero-event bound.}
With beta quantile $B^{-1}$, inversion of the binomial tails
$\Pr_L(F\ge f)=0.025$, $\Pr_U(F\le f)=0.025$ gives the pointwise 95\%
Clopper--Pearson interval \citep{clopper1934limits}, including endpoint cases:
\begin{align}
 L&=\begin{cases}0,&f=0,\\
 B^{-1}(0.025;f,M_c-f+1),&f>0,\end{cases}\\
 U&=\begin{cases}1,&f=M_c,\\
 B^{-1}(0.975;f+1,M_c-f),&f<M_c.\end{cases}
 \label{eq:r33-cp}
\end{align}
For $c$ perfect payloads among $R$ records, substitute $f=c,M_c=R$;
the event and independent unit change. SciPy beta quantiles are checked
against its exact binomial-test interval \citep{scipy1153binomialci}.
These exact intervals can be conservative \citep{thulin2014cost}.
For zero observed failures, $(1-q)^{M_c}=0.05$ gives the
\emph{one-sided} 95\% bound
\begin{equation}
 q_U=1-0.05^{1/M_c};
 \qquad M_c=4480000\ \Longrightarrow\
 q_U\simeq6.6869\times10^{-7}.
 \label{eq:r33-zero}
\end{equation}
Zero observations do not establish a zero underlying rate.

\paragraph{Planning a new record from a nested design bank.}
R3.4 uses five earlier Gaussian banks with eight BSC masks each.
Conditional on Gaussian input $Z$, write $\mu(Z)=\mathbb E[d\mid Z]$ and
$v(Z)=\operatorname{Var}(d\mid Z)$. Total variance gives
$\operatorname{Var}(d)=\mathbb E[v(Z)]+\operatorname{Var}[\mu(Z)]$;
bank-mean variance estimates the second term plus $\mathbb E[v(Z)]/8$.
With within-bank sample variance $s_g^2$ and mean $\overline d_g$, subtract
that reduced contribution (truncate at zero) and add the full within-bank variance:
\begin{equation}
 \widehat V=
 \overline{s_g^2}+
 \max\!\left(s_{\overline d_g}^2-\frac{\overline{s_g^2}}8,0\right).
 \label{eq:r34-planning}
\end{equation}

\paragraph{Contrasts and interaction bounds on a joint event.}
On the joint primary-interval event, conservative contrasts satisfy
\begin{equation}
 \mu_a-\mu_b\ \in\ [L_a-U_b,\ U_a-L_b].
 \label{eq:r43-contrast}
\end{equation}
For $I=\mu_{00}-\mu_{10}-\mu_{01}+\mu_{11}$,
\begin{align}
 L_I&=L_{00}-U_{10}-U_{01}+L_{11},\nonumber\\
 U_I&=U_{00}-L_{10}-L_{01}+U_{11}.
 \label{eq:r43-interaction-ci}
\end{align}

For paired $\delta_r=d_{a,r}-d_{b,r}$, expanding centred squares gives
\begin{equation}
 s_\delta^2=s_a^2+s_b^2-2s_{ab},\qquad
 s_{ab}=\frac1{R-1}\sum_r(d_{a,r}-\bar d_a)(d_{b,r}-\bar d_b).
 \label{eq:r43-paired-ci}
\end{equation}
Apply Eq.~\eqref{eq:app-mean-interval} to $\delta$ with $R=1400,m=1$
for pointwise paired intervals, distinct from the joint primary bounds.

\paragraph{Bundle-level paired objectives.}
The 12 unique main conditions have equal design weights, not inferred
channel occurrence probabilities. Define the objective and paired gain once:
\begin{align}
 J_{fc,r}&=\frac1{12}\sum_{s\in\mathcal S_{\rm main}}d_{fc,s,r},\quad
 \bar J_{fc}=R_v^{-1}\sum_rJ_{fc,r},\label{eq:r42-objective}\\
 g_{v,r}&=J_{{\rm base},r}-J_{v,r}
 =\frac1{12}\sum_{s\in\mathcal S_{\rm main}}(d_{{\rm base},s,r}-d_{v,s,r}).
 \label{eq:r51-gain}
\end{align}

\paragraph{The R5.2 joint constraint.}
At the prespecified anchor-point scenario, the upper mean bound from
Eq.~\eqref{eq:app-mean-interval}, with $R=350,m=19$, must satisfy
\begin{equation}
 U_v\leq\tau,\qquad r_{t,v}\leq320000\ {\rm bit/s}.
 \label{eq:r52-quality}
\end{equation}
The quality threshold is unchanged; this is a population-mean decision,
not a per-record guarantee.


\paragraph{One pilot precision-planning relation.}
Using phase-0 standard deviations $\widetilde s_j$, choose the smallest
allowed count meeting every prespecified half-width target:
\begin{equation}
 t_{1-0.05/(2m),R-1}\frac{\widetilde s_j}{\sqrt R}\leq\epsilon_j
 \quad\text{for all }j,\qquad R\geq200,\quad R\in50\mathbb N.
 \label{eq:app-pilot-planning}
\end{equation}
R4 uses $m=92$, $\epsilon_j=0.05\widetilde\mu_j$ and selects 1400.
R5 uses $m=19$, $\epsilon_j=\max(0.01,0.05|\widetilde\mu_j|)$ and selects
350; 300 fails some pilot targets. R3.4 uses $m=30$, $\epsilon=0.1\tau$,
$\widetilde s_j=\sqrt{\widehat V_j}$ for uncoded $p=0.0001$ and coded
$p=0.003$, selecting 300. Planning does not guarantee realised precision:
7/92 R4 means miss the 5\% relative half-width target, with maximum 42.43\%.
This is uncertainty in estimated means, not reconstruction error.
No confirmation data enter planning or trigger extra repeats.

\subsection{Parameter-mapping convention}
\label{app:random}
The declared project mapping $p_{\rm ID}=10^{-4+3u}$, uniform $u\in[0,1)$,
gives CDF $(\log_{10}p+4)/3$. Differentiation yields
\begin{equation}
 f_{p_{\rm ID}}(p)=\frac{1}{3\ln(10)\,p},
 \qquad 10^{-4}\le p<0.1.
 \label{eq:r33-density}
\end{equation}
Equal logarithmic intervals receive equal weight; this is a documented
project choice, not a supplied external rule.
